\documentclass[aspectratio=169,11pt]{beamer}

\providecommand\UseMathForPositioningText{}

% ============================================================
% بسته‌های پایه
% ============================================================
\usepackage{amsmath,amssymb,amsthm}
\usepackage{mathrsfs}
\usepackage{booktabs}
\usepackage{array}
\usepackage{xcolor}
\usepackage{tikz}
\usetikzlibrary{positioning, arrows.meta, shapes.geometric, calc, shadows}
\usepackage{fontawesome5}

% ============================================================
% بسته فارسی
% ============================================================
\usepackage{xepersian}

\settextfont{Vazirmatn}
% اگر Vazirmatn نداری، یکی از این‌ها را فعال کن:
% \settextfont{XB Niloofar}
% \settextfont{Yas}
% \settextfont{XB Zar}
% \settextfont{Amiri}
% \settextfont{FreeSerif}

% ============================================================
% قالب
% ============================================================
\usetheme{metropolis}

% ============================================================
% پالت رنگی
% ============================================================
\definecolor{MainBlue}{RGB}{26, 54, 93}
\definecolor{AccentOrange}{RGB}{230, 126, 34}
\definecolor{SoftGray}{RGB}{245, 246, 250}
\definecolor{TextDark}{RGB}{45, 55, 72}
\definecolor{SuccessGreen}{RGB}{39, 174, 96}
\definecolor{AlertRed}{RGB}{192, 57, 43}

\setbeamercolor{normal text}{fg=TextDark, bg=white}
\setbeamercolor{alerted text}{fg=AlertRed}
\setbeamercolor{example text}{fg=SuccessGreen}
\setbeamercolor{palette primary}{fg=white, bg=MainBlue}
\setbeamercolor{frametitle}{fg=white, bg=MainBlue}
\setbeamercolor{title}{fg=MainBlue}
\setbeamercolor{structure}{fg=MainBlue}
\setbeamercolor{block title}{fg=white, bg=MainBlue}
\setbeamercolor{block body}{fg=TextDark, bg=SoftGray}
\setbeamercolor{block title alerted}{fg=white, bg=AlertRed}
\setbeamercolor{block body alerted}{fg=TextDark, bg=SoftGray}
\setbeamercolor{block title example}{fg=white, bg=SuccessGreen}
\setbeamercolor{block body example}{fg=TextDark, bg=SoftGray}
\setbeamercolor{section in toc}{fg=MainBlue}
\setbeamercolor{subsection in toc}{fg=AccentOrange}

% ============================================================
% تنظیمات ظاهری
% ============================================================
\setbeamertemplate{navigation symbols}{}
\setbeamertemplate{footline}[frame number]
\setbeamerfont{frametitle}{size=\large, series=\bfseries}
\setbeamerfont{title}{size=\LARGE, series=\bfseries}
\setbeamerfont{subtitle}{size=\large}
\setbeamerfont{section in toc}{size=\small, series=\bfseries}
\setbeamertemplate{blocks}[rounded][shadow=false]

\setbeamertemplate{itemize item}{\textcolor{AccentOrange}{\small$\blacktriangleright$}}
\setbeamertemplate{itemize subitem}{\textcolor{MainBlue}{\tiny$\bullet$}}
\setbeamertemplate{enumerate item}{\textcolor{AccentOrange}{\textbf{\arabic{enumi}.}}}

\definecolor{MainBlue}{RGB}{26, 54, 93}
\definecolor{DeepBlue}{RGB}{15, 35, 65}
\definecolor{AccentOrange}{RGB}{230, 126, 34}
\definecolor{SoftGray}{RGB}{245, 246, 250}
\definecolor{TextDark}{RGB}{45, 55, 72}
\definecolor{SuccessGreen}{RGB}{39, 174, 96}
\definecolor{AlertRed}{RGB}{192, 57, 43}
% ============================================================
% استایل‌های TikZ
% ============================================================
\tikzset{
  flowchartbox/.style={
    rectangle, rounded corners, draw=MainBlue, fill=SoftGray,
    text width=10cm, align=center, minimum height=0.6cm,
    font=\footnotesize, inner sep=5pt
  },
  flowchartarrow/.style={
    -{Stealth[length=2.5mm]}, thick, AccentOrange
  },
  stepbox/.style={
    rectangle, rounded corners, draw=MainBlue, fill=SoftGray,
    text width=11cm, align=center, minimum height=0.55cm,
    font=\scriptsize, inner sep=4pt
  }
}

% ============================================================
% اطلاعات عنوان
% ============================================================
\title{انتظارات عقلایی و نظریه سیاست اقتصادی}
\subtitle{\lr{Sargent, T. \& Wallace, N. (1981)}}
\author{توماس سارجنت و نیل والاس}
\institute{}
\date{}

% ============================================================
% دستورات سفارشی
% ============================================================
\newcommand{\highlight}[1]{\textcolor{AccentOrange}{\textbf{#1}}}
\newcommand{\important}[1]{\textcolor{AlertRed}{\textbf{#1}}}
\newcommand{\keyword}[1]{\textcolor{MainBlue}{\textbf{#1}}}

\begin{document}

% ============================================================
% اسلاید عنوان زیبا
% ============================================================
{
\setbeamertemplate{background}{
	\begin{tikzpicture}[remember picture, overlay]
		\fill[MainBlue] (current page.south west) rectangle (current page.north east);
		\fill[AccentOrange, opacity=0.15]
		(current page.south west) rectangle ([yshift=0.35\paperheight]current page.north west);
		\node[opacity=0.05, text=white, font=\fontsize{180}{180}\selectfont]
		at (current page.center) {\faChartLine};
	\end{tikzpicture}
}
\begin{frame}[plain]
	\begin{center}
		\vspace{0.3cm}

		{\color{white}\LARGE\bfseries انتظارات عقلایی و نظریه سیاست اقتصادی}\\[0.35cm]

		{\color{AccentOrange}\normalsize \lr{Sargent, T. \& Wallace, N. (1981)}}\\[0.5cm]

		{\color{AccentOrange}\rule{6cm}{0.6pt}}\\[0.5cm]

		{\color{white}\small ارائه‌دهنده}\\[0.1cm]
		{\color{AccentOrange}\large\bfseries مسعود صنعتگر قوچانی}\\[0.4cm]

		{\color{white}\small استاد درس}\\[0.1cm]
		{\color{white}\large\bfseries دکتر مسعود محمدی الموتی}\\[0.4cm]

		{\color{white}\small تاریخ ارائه}\\[0.1cm]
		{\color{AccentOrange}\normalsize ۲۱ مهر ۱۴۰۵}

	\end{center}
\end{frame}
}% فهرست مطالب (بخش اول راست، بخش دوم چپ)
% ============================================================
\begin{frame}{فهرست مطالب}
	\begin{LTR}
		\begin{columns}[t]
			\begin{column}{0.48\textwidth}
				\begin{RTL}
					\begin{block}{\faListUl\ بخش دوم}
						\tableofcontents[sections={5-8}]
					\end{block}
				\end{RTL}
			\end{column}
			\hfill
			\begin{column}{0.48\textwidth}
				\begin{RTL}
					\begin{block}{\faListUl\ بخش اول}
						\tableofcontents[sections={1-4}]
					\end{block}
				\end{RTL}
			\end{column}
		\end{columns}
	\end{LTR}
\end{frame}

% ============================================================
\section{معرفی مقاله}
% ============================================================

\begin{frame}{اطلاعات کتاب‌شناختی}
	\begin{LTR}
		\begin{columns}[T]
			\begin{column}{0.60\textwidth}
				\begin{RTL}
					\begin{itemize}
						\item \keyword{نویسندگان:} \lr{Thomas J. Sargent} و \lr{Neil Wallace}
						\item \keyword{عنوان:} \lr{Rational Expectations and the Theory of Economic Policy}
						\item \keyword{منبع:} فصل ۳۱ در کتاب
						\item \keyword{انتشار اولیه:} \lr{Working Paper No. 29} (۱۹۷۴)
						\item \keyword{انتشار رسمی:} \lr{Journal of Monetary Economics}، ۱۹۷۶
					\end{itemize}
				\end{RTL}
			\end{column}
			\hfill
			\begin{column}{0.35\textwidth}
				\begin{RTL}
					\begin{block}{}
						\centering
						{\color{MainBlue}\Huge\faBookOpen}\\[0.4cm]
						{\footnotesize\lr{University of Minnesota Press}}\\[0.2cm]
						{\small ۱۹۸۱}
					\end{block}
				\end{RTL}
			\end{column}
		\end{columns}
	\end{LTR}
\end{frame}

\begin{frame}{جایگاه تاریخی مقاله}
	\begin{block}{سه ستون انقلاب انتظارات عقلایی}
		\begin{enumerate}
			\item \lr{Lucas (1972): Expectations and the Neutrality of Money}
			\item \highlight{\lr{Sargent \& Wallace (1976/1981)}:} قضیه بی‌اثری سیاست
			\item \lr{Lucas (1976): Econometric Policy Evaluation: A Critique}
		\end{enumerate}
	\end{block}
	\vspace{0.4cm}
	\begin{LTR}
		\begin{columns}[T]
			\begin{column}{0.48\textwidth}
				\begin{RTL}
					\begin{alertblock}{زمینه تاریخی}
						\centering
						رکود تورمی دهه ۱۹۷۰
					\end{alertblock}
				\end{RTL}
			\end{column}
			\hfill
			\begin{column}{0.48\textwidth}
				\begin{RTL}
					\begin{exampleblock}{پیامد}
						\centering
						چالش با سیاست‌های تثبیتی کینزی
					\end{exampleblock}
				\end{RTL}
			\end{column}
		\end{columns}
	\end{LTR}
\end{frame}

\begin{frame}{پرسش اصلی مقاله}
	\begin{block}{پرسش محوری}
		\centering
		\large
		آیا دولت می‌تواند با استفاده از \highlight{سیاست‌های پولی نظام‌مند}، متغیرهای واقعی اقتصاد مانند تولید و بیکاری را به‌طور دائمی تحت تأثیر قرار دهد؟
	\end{block}
	\vspace{0.5cm}
	\begin{alertblock}{پاسخ مقاله}
		\centering
		\large
		\important{خیر.} تحت فرض انتظارات عقلایی، هر سیاست نظام‌مند و قابل پیش‌بینی، بی‌اثر است.
		\vspace{0.2cm}

		این همان \important{قضیه بی‌اثری سیاست} است.
	\end{alertblock}
\end{frame}

\begin{frame}{ساختار منطقی مقاله}
	\begin{center}
		\begin{tikzpicture}[node distance=0.6cm]
			\node[flowchartbox] (a) {\begin{RTL}\textbf{فرض:} عوامل اقتصادی انتظارات عقلایی دارند\end{RTL}};
			\node[flowchartbox, below=of a] (b) {\begin{RTL}\textbf{نتیجه ۱:} تنها خطاهای پیش‌بینی بر تولید اثر دارند\end{RTL}};
			\node[flowchartbox, below=of b] (c) {\begin{RTL}\textbf{نتیجه ۲:} سیاست نظام‌مند قابل پیش‌بینی، خطای پیش‌بینی ایجاد نمی‌کند\end{RTL}};
			\node[flowchartbox, below=of c, fill=AccentOrange!20, draw=AccentOrange] (d) {\begin{RTL}\textbf{نتیجه نهایی:} سیاست نظام‌مند بی‌اثر است\end{RTL}};

			\draw[flowchartarrow] (a) -- (b);
			\draw[flowchartarrow] (b) -- (c);
			\draw[flowchartarrow] (c) -- (d);
		\end{tikzpicture}
	\end{center}
\end{frame}

% ============================================================
\section{فرض کلیدی: انتظارات عقلایی}
% ============================================================

\begin{frame}{انتظارات عقلایی: تعریف}
	\begin{block}{تعریف جان ماث (\lr{John Muth, 1961})}
		انتظارات عقلایی یعنی انتظارات ذهنی عوامل اقتصادی با \highlight{پیش‌بینی ریاضی بهینه} منطبق است.
	\end{block}
	\vspace{0.3cm}
	\begin{LTR}
		\begin{columns}[T]
			\begin{column}{0.52\textwidth}
				\begin{RTL}
					\begin{itemize}
						\item \faCheckCircle\ مردم \keyword{همه اطلاعات} را پردازش می‌کنند
						\item \faCheckCircle\ \keyword{ساختار مدل} را می‌شناسند
						\item \faCheckCircle\ \keyword{قواعد سیاستی} را می‌شناسند
					\end{itemize}
				\end{RTL}
			\end{column}
			\hfill
			\begin{column}{0.44\textwidth}
				\begin{RTL}
					\begin{exampleblock}{نتیجه}
						انتظارات \highlight{خطای سیستماتیک} ندارد
					\end{exampleblock}
				\end{RTL}
			\end{column}
		\end{columns}
	\end{LTR}
\end{frame}

\begin{frame}{تفاوت با انتظارات تطبیقی}
	\begin{center}
		\renewcommand{\arraystretch}{1.4}
		\begin{tabular}{ccc}
			\toprule
			\textbf{ویژگی} & \textbf{\textcolor{AlertRed}{انتظارات تطبیقی}} & \textbf{\textcolor{SuccessGreen}{انتظارات عقلایی}} \\
			\midrule
			جهت‌گیری        & گذشته‌نگر                                       & آینده‌نگر                                           \\
			اطلاعات        & فقط گذشته                                      & همه اطلاعات موجود                                  \\
			خطای سیستماتیک & دارد                                           & ندارد                                              \\
			سرعت یادگیری   & کند                                            & فوری                                               \\
			مبنای نظری     & تجربی                                          & بهینه‌سازی                                          \\
			\bottomrule
		\end{tabular}
	\end{center}
\end{frame}

\begin{frame}{پیامد انتظارات عقلایی برای سیاست}
	\begin{center}
		\renewcommand{\arraystretch}{1.5}
		\begin{tabular}{lcc}
			\toprule
			\textbf{نوع سیاست} & \textbf{آیا پیش‌بینی می‌شود؟} & \textbf{اثر بر تولید} \\
			\midrule
			قاعده نظام‌مند      & \faCheckCircle              & \important{صفر}       \\
			شوک تصادفی         & \faTimesCircle              & موقتی و بی‌فایده       \\
			\bottomrule
		\end{tabular}
	\end{center}
	\vspace{0.5cm}
	\begin{alertblock}{نکته کلیدی}
		تنها بخش \important{غیرقابل پیش‌بینی} سیاست می‌تواند اثر واقعی داشته باشد.
	\end{alertblock}
\end{frame}

% ============================================================
\section{ساختار مدل ریاضی}
% ============================================================

\begin{frame}{معادله عرضه کل}
	\begin{block}{منحنی فیلیپس با انتظارات عقلایی}
		\centering
		\large
		\[
			y_t = \bar{y} + \alpha (p_t - E_{t-1}p_t) + \varepsilon_t
		\]
	\end{block}
	\vspace{0.2cm}
	\begin{LTR}
		\begin{columns}[T]
			\begin{column}{0.48\textwidth}
				\begin{RTL}
					\begin{itemize}
						\item $y_t$: تولید واقعی
						\item $\bar{y}$: تولید طبیعی
						\item $p_t$: سطح قیمت
					\end{itemize}
				\end{RTL}
			\end{column}
			\hfill
			\begin{column}{0.48\textwidth}
				\begin{RTL}
					\begin{itemize}
						\item $E_{t-1}p_t$: انتظار قیمت
						\item $\varepsilon_t$: شوک عرضه
						\item $\alpha > 0$: حساسیت
					\end{itemize}
				\end{RTL}
			\end{column}
		\end{columns}
	\end{LTR}
\end{frame}

\begin{frame}{تفسیر اقتصادی معادله عرضه}
	\begin{center}
		\renewcommand{\arraystretch}{1.5}
		\begin{tabular}{lc}
			\toprule
			\textbf{وضعیت}     & \textbf{اثر بر تولید}                  \\
			\midrule
			$p_t > E_{t-1}p_t$ & \textcolor{SuccessGreen}{افزایش تولید} \\
			$p_t = E_{t-1}p_t$ & تولید در سطح طبیعی                     \\
			$p_t < E_{t-1}p_t$ & \textcolor{AlertRed}{کاهش تولید}       \\
			\bottomrule
		\end{tabular}
	\end{center}
	\vspace{0.5cm}
	\begin{alertblock}{نکته کلیدی}
		فقط \important{خطای پیش‌بینی قیمت} اثر واقعی دارد.
	\end{alertblock}
\end{frame}

\begin{frame}{معادله تقاضای کل}
	\begin{block}{معادله مبادله‌ای}
		\centering
		\large
		\[
			m_t + v_t = p_t + y_t
		\]
	\end{block}
	\vspace{0.3cm}
	\begin{LTR}
		\begin{columns}[T]
			\begin{column}{0.48\textwidth}
				\begin{RTL}
					\begin{itemize}
						\item $m_t$: عرضه پول
						\item $v_t$: شوک سرعت گردش
					\end{itemize}
				\end{RTL}
			\end{column}
			\hfill
			\begin{column}{0.48\textwidth}
				\begin{RTL}
					\begin{itemize}
						\item $p_t$: سطح قیمت
						\item $y_t$: تولید واقعی
					\end{itemize}
				\end{RTL}
			\end{column}
		\end{columns}
	\end{LTR}
	\vspace{0.3cm}
	\begin{exampleblock}{حل برای قیمت}
		\centering
		$p_t = m_t + v_t - y_t$
	\end{exampleblock}
\end{frame}

\begin{frame}{قاعده سیاست پولی}
	\begin{block}{قاعده بازخوردی دولت}
		\centering
		\large
		\[
			m_t = \mu_0 + \mu_1 y_{t-1} + \mu_2 m_{t-1} + u_t
		\]
	\end{block}
	\vspace{0.3cm}
	\begin{itemize}
		\item $\mu_0, \mu_1, \mu_2$: \keyword{پارامترهای سیاست}
		\item $u_t$: \highlight{شوک غیرقابل پیش‌بینی} سیاست
	\end{itemize}
	\vspace{0.3cm}
	\begin{alertblock}{هدف تحلیل}
		آیا تغییر $\mu_1$ می‌تواند تولید واقعی را تحت تأثیر قرار دهد؟
	\end{alertblock}
\end{frame}

% ============================================================
\section{حل کامل مدل}
% ============================================================

\begin{frame}{مرحله ۱: قیمت تعادلی}
	با جایگذاری معادله عرضه در تقاضا:
	\[
		p_t = m_t + v_t - \bar{y} - \alpha(p_t - E_{t-1}p_t) - \varepsilon_t
	\]
	\vspace{0.3cm}
	با مرتب کردن:
	\[
		p_t (1 + \alpha) = m_t + v_t - \bar{y} + \alpha E_{t-1}p_t - \varepsilon_t
	\]
	\vspace{0.3cm}
	\begin{block}{قیمت تعادلی}
		\centering
		\[
			p_t = \frac{1}{1+\alpha}\left(m_t + v_t - \bar{y} + \alpha E_{t-1}p_t - \varepsilon_t\right)
		\]
	\end{block}
\end{frame}

\begin{frame}{مرحله ۲: محاسبه انتظار قیمت}
	\begin{block}{انتظار عقلایی قیمت}
		\[
			E_{t-1}p_t = \frac{1}{1+\alpha}\left(E_{t-1}m_t + E_{t-1}v_t - \bar{y} + \alpha E_{t-1}p_t - E_{t-1}\varepsilon_t\right)
		\]
	\end{block}
	\vspace{0.3cm}
	با فرض:
	\[
		E_{t-1}v_t = 0, \quad E_{t-1}\varepsilon_t = 0
	\]
	و از قاعده سیاست:
	\[
		E_{t-1}m_t = \mu_0 + \mu_1 y_{t-1} + \mu_2 m_{t-1}
	\]
\end{frame}

\begin{frame}{مرحله ۳: خطای پیش‌بینی قیمت}
	\begin{block}{تفریق دو معادله}
		\[
			p_t - E_{t-1}p_t = \frac{1}{1+\alpha}\left[(m_t - E_{t-1}m_t) + v_t - \varepsilon_t\right]
		\]
	\end{block}
	\vspace{0.3cm}
	و چون:
	\[
		m_t - E_{t-1}m_t = u_t
	\]
	\vspace{0.3cm}
	\begin{alertblock}{خطای پیش‌بینی قیمت}
		\centering
		\[
			p_t - E_{t-1}p_t = \frac{1}{1+\alpha}(u_t + v_t - \varepsilon_t)
		\]
	\end{alertblock}
\end{frame}

\begin{frame}{مرحله ۴: معادله نهایی تولید}
	\[
		y_t = \bar{y} + \alpha \cdot \frac{1}{1+\alpha}(u_t + v_t - \varepsilon_t) + \varepsilon_t
	\]
	\vspace{0.5cm}
	\begin{block}{معادله نهایی}
		\centering
		\Large
		\[
			\boxed{y_t = \bar{y} + \frac{\alpha}{1+\alpha}u_t + \frac{\alpha}{1+\alpha}v_t + \frac{1}{1+\alpha}\varepsilon_t}
		\]
	\end{block}
\end{frame}

% ============================================================
\section{قضیه بی‌اثری سیاست}
% ============================================================

\begin{frame}{حذف پارامترهای سیاست}
	\begin{block}{معادله نهایی تولید}
		\[
			y_t = \bar{y} + \frac{\alpha}{1+\alpha}u_t + \frac{\alpha}{1+\alpha}v_t + \frac{1}{1+\alpha}\varepsilon_t
		\]
	\end{block}
	\vspace{0.3cm}
	\begin{alertblock}{مشاهده کلیدی}
		\centering
		\large
		پارامترهای $\mu_0, \mu_1, \mu_2$ به‌طور کامل \important{حذف شده‌اند!}
	\end{alertblock}
	\vspace{0.3cm}
	\begin{itemize}
		\item توزیع $y_t$ \keyword{هیچ وابستگی} به قاعده سیاستی ندارد
		\item تغییر $\mu_1$ \important{هیچ اثری} بر تولید ندارد
		\item تنها \highlight{شوک‌های غیرقابل پیش‌بینی} اثر دارند
	\end{itemize}
\end{frame}

\begin{frame}{تحلیل واریانس تولید}
	\begin{block}{واریانس تولید}
		\[
			\text{Var}(y_t) = \left(\frac{\alpha}{1+\alpha}\right)^2 (\sigma_u^2 + \sigma_v^2) + \left(\frac{1}{1+\alpha}\right)^2 \sigma_\varepsilon^2
		\]
	\end{block}
	\vspace{0.3cm}
	\begin{LTR}
		\begin{columns}[T]
			\begin{column}{0.32\textwidth}
				\begin{RTL}
					\centering
					$\sigma_u^2$\\[0.2cm]
					\small شوک سیاست
				\end{RTL}
			\end{column}
			\hfill
			\begin{column}{0.32\textwidth}
				\begin{RTL}
					\centering
					$\sigma_v^2$\\[0.2cm]
					\small شوک سرعت
				\end{RTL}
			\end{column}
			\hfill
			\begin{column}{0.32\textwidth}
				\begin{RTL}
					\centering
					$\sigma_\varepsilon^2$\\[0.2cm]
					\small شوک عرضه
				\end{RTL}
			\end{column}
		\end{columns}
	\end{LTR}
	\vspace{0.3cm}
	\begin{alertblock}{نکته}
		هیچ‌کدام از پارامترهای سیاست در واریانس ظاهر نمی‌شوند.
	\end{alertblock}
\end{frame}

\begin{frame}{تفسیر شهودی}
	\begin{center}
		\begin{tikzpicture}[node distance=0.35cm]
			\node[stepbox] (s1) {\begin{RTL}\textbf{۱.} مردم قاعده سیاستی $\mu_1$ را می‌دانند\end{RTL}};
			\node[stepbox, below=of s1] (s2) {\begin{RTL}\textbf{۲.} آن را در انتظارات قیمتی $E_{t-1}p_t$ وارد می‌کنند\end{RTL}};
			\node[stepbox, below=of s2] (s3) {\begin{RTL}\textbf{۳.} دولت پول را افزایش می‌دهد، اما مردم قبلاً پیش‌بینی کرده‌اند\end{RTL}};
			\node[stepbox, below=of s3] (s4) {\begin{RTL}\textbf{۴.} خطای پیش‌بینی $p_t - E_{t-1}p_t$ تغییر نمی‌کند\end{RTL}};
			\node[stepbox, below=of s4, fill=AccentOrange!20, draw=AccentOrange] (s5) {\begin{RTL}\textbf{۵.} تولید واقعی تکان نمی‌خورد\end{RTL}};

			\foreach \i/\j in {s1/s2, s2/s3, s3/s4, s4/s5}
			\draw[flowchartarrow] (\i) -- (\j);
		\end{tikzpicture}
	\end{center}
\end{frame}

\begin{frame}{پیام‌های سیاستی}
	\begin{LTR}
		\begin{columns}[T]
			\begin{column}{0.48\textwidth}
				\begin{RTL}
					\begin{alertblock}{اگر دولت...}
						\begin{itemize}
							\item قاعده پولی را تغییر دهد
							\item سیاست را نظام‌مند کند
							\item شوک وارد کند
						\end{itemize}
					\end{alertblock}
				\end{RTL}
			\end{column}
			\hfill
			\begin{column}{0.48\textwidth}
				\begin{RTL}
					\begin{exampleblock}{آنگاه...}
						\begin{itemize}
							\item تولید تغییر نمی‌کند
							\item فقط تورم ایجاد می‌شود
							\item نوسان موقتی رخ می‌دهد
						\end{itemize}
					\end{exampleblock}
				\end{RTL}
			\end{column}
		\end{columns}
	\end{LTR}
	\vspace{0.4cm}
	\begin{block}{پیام اصلی}
		\centering
		\large
		\important{سیاست تثبیتی نظام‌مند بی‌اثر است.}
	\end{block}
\end{frame}

% ============================================================
\section{اهمیت و نقدها}
% ============================================================

\begin{frame}{اهمیت و جایگاه}
	\begin{block}{میراث مقاله}
		\begin{itemize}
			\item ستون \keyword{مکتب کلاسیک جدید}
			\item همراه با \highlight{نقد لوکاس} (۱۹۷۶)
			\item چالشی بنیادین با \important{سیاست‌های تثبیتی کینزی}
		\end{itemize}
	\end{block}
	\vspace{0.3cm}
	\begin{exampleblock}{زمینه‌ساز تحولات بعدی}
		\begin{itemize}
			\item هدف‌گذاری تورم
			\item استقلال بانک مرکزی
			\item قاعده تیلور
		\end{itemize}
	\end{exampleblock}
\end{frame}

\begin{frame}{نقدهای وارد بر مقاله}
	\begin{LTR}
		\begin{columns}[T]
			\begin{column}{0.48\textwidth}
				\begin{RTL}
					\begin{alertblock}{\faExclamationTriangle\ نقدهای نظری}
						\begin{itemize}
							\item چسبندگی قیمت و دستمزد
							\item اطلاعات نامتقارن
							\item مدل‌های کینزی جدید
						\end{itemize}
					\end{alertblock}
				\end{RTL}
			\end{column}
			\hfill
			\begin{column}{0.48\textwidth}
				\begin{RTL}
					\begin{alertblock}{\faExclamationTriangle\ نقدهای تجربی}
						\begin{itemize}
							\item شواهد اثر سیاست نظام‌مند
							\item نوسانات واقعی اقتصاد
							\item نقش انتظارات غیرعقلایی
						\end{itemize}
					\end{alertblock}
				\end{RTL}
			\end{column}
		\end{columns}
	\end{LTR}
	\vspace{0.3cm}
	\begin{block}{جمع‌بندی}
		\centering
		با وجود نقدها، این مقاله نقطه عطفی در تاریخ اقتصاد کلان است.
	\end{block}
\end{frame}

% ============================================================
\section{جمع‌بندی}
% ============================================================

\begin{frame}{جمع‌بندی نهایی}
	\begin{block}{پیام اصلی}
		\centering
		\large
		تحت فرض انتظارات عقلایی، \important{سیاست پولی نظام‌مند نمی‌تواند متغیرهای واقعی را تحت تأثیر قرار دهد.}
	\end{block}
	\vspace{0.5cm}
	\begin{LTR}
		\begin{columns}[T]
			\begin{column}{0.48\textwidth}
				\begin{RTL}
					\begin{itemize}
						\item \keyword{مکانیسم:} مردم قاعده را می‌دانند
						\item \keyword{نتیجه:} پارامترها حذف می‌شوند
					\end{itemize}
				\end{RTL}
			\end{column}
			\hfill
			\begin{column}{0.48\textwidth}
				\begin{RTL}
					\begin{itemize}
						\item \keyword{پیام:} فقط شوک‌های غیرمنتظره اثر دارند
						\item \keyword{میراث:} کلاسیک جدید
					\end{itemize}
				\end{RTL}
			\end{column}
		\end{columns}
	\end{LTR}
\end{frame}

\begin{frame}[standout]
	\centering
	\Huge
	\textcolor{white}{پایان مقاله اول}
	\vspace{0.5cm}

	\large
	\textcolor{white}{آماده برای مقاله دوم:}\\
	\textcolor{AccentOrange}{\lr{\textbf{After Keynesian Macroeconomics}}}
\end{frame}

\end{document}
